This bachelor's thesis investigates how to build and evaluate an Estonian wake-word detector for a resource-constrained smart-home microcontroller. The work focuses on the \enquote{Kratt} project, whose target is local detection of the phrase \enquote{Kuule Kratt} on an ESP32-S3 class device without relying on cloud services.

The thesis reconstructs a \texttt{microWakeWord}-based training and evaluation pipeline and validates it with a public control experiment on the \texttt{Speech Commands} dataset. The initial clip-level evaluation proved misleading due to data leakage, so independent held-out test sets were constructed, train--test disjointness checks were added, and streaming-mode FAPH (false accepts per hour) was adopted as the central metric.

The results show that no current model or combination of models satisfies all targets simultaneously: low FAPH, high recall on real speakers, and accurate phrase distinction. The main contribution is therefore a reproducible evaluation protocol and an empirical description of the trade-offs that arise in small-resource-language local wake-word development, rather than a final deployable model.


% No need to change this
The thesis is in \langEng~and contains \calculatepages pages of text, 
\total{totalchapters} chapters\ifthenelse{\equal{\totvalue{figure}}{0}}{}{% If no figures, do nothing
, \total{figure} \ifnum\totvalue{figure}=1 figure\else figures\fi%
}\ifthenelse{\equal{\totvalue{table}}{0}}{}{% If no tables, do nothing
, \total{table} \ifnum\totvalue{table}=1 table\else tables\fi%
}.
